\section[Pfister-Roumy]{Création de la liste Pfister-Roumy}

Ancien pôle de référence nationale sous le statut de \enquote{Centre d'acquisition et de diffusion de l'information scientifique et technique} (CADIST), la bibliothèque Cujas présente une trajectoire singulière. En effet, alors que la majorité des grandes bibliothèques ont constitué leurs fonds patrimoniaux à partir de confiscations révolutionnaires. Sa collection d'origine comprend celle \enquote{de la Faculté de Droit de Paris [qui] a été dispersée sous la Révolution}. Le fonds actuel a été reconstruit \enquote{rétropectivement par des bibliothécaires, dans le but de reconstituer \enquote{le cabinet de travail d'un robin d'Ancien Régime} -- telle était, du moins, la doxa transmise jusqu'à présent de bouche à oreille de bibliothécaire}. Ce fonds patrimonial s'est essentiellement constitué à partir des années 1870, grâce à des achats, dons ou legs. 

La constitution de ce fonds comprend : \enquote{12 000 volumes imprimés, datant de la fin du XVe siècle au début du XIXe siècle, essentiellement des livres des XVIIe et XVIIIe siècles, 7 incunables, quelques ouvrages précieux postérieurs à 1800 et quelques périodiques. Le quart de ce fonds provient de pays étrangers : Allemagne, Pays-Bas, Angleterre et Italie. Particulièrement riche, il comprend 25 \% d'ouvrages non référencés à la BnF. La réserve conserve également des cours et notes de cours manuscrits, ainsi qu'une partie des archives de la Faculté de Droit.}. Ces fonds patrimoniaux représentent ainsi les grands courants du droit : droit civil, criminel et canonique, pratique du droit ancien ou nouveau (après 1789), coutumiers, oeuvres des grands jurisconsultes français et étrangers.
	
Afin de valoriser la richesse de ce patrimoine, la bibliothèque Cujas conçoit le projet de développer sa propre bibliothèque, une ambition émergeant en même temps que la numérisation des cours de Jean Carbonnier. A l'origine, l''établissement souhaite s'insérer dans une dynamique nationale en postulant à l'appel à projet de Persée, dans l'espoir de devenir un pôle de réplication. C'est dans cette perspective que la bibliothèque acquiert un scanner Kirtas destiné à industrialiser sa chaîne de traitement interne. Bien que cette candidature n'ait finalement pas été retenue, elle sera le point de départ de réflexion de la bibliothèque souhaitant capitaliser sur cette  acquistion. 

Ne souhaitant pas le laisser inutilisée, la direction de la bibliothèque Cujas opère un changement stratégique majeur. Elle choisit de s'associer étroitement avec deux éminents  historiens du droit de la Faculté de droit de Paris, les professeurs Laurent Pfister et Franck Roumy. Cette collaboration inédite entre les professionnels de bibliothèques et les enseignants-chercheurs permet d'élaborer une première liste de référence scientifique, désignée plus tard sous le nom de \enquote{liste Pfister-Roumy}. Contrairement au projet Carbonnier, ce nouveau projet ne se concentre plus sur des cours polycopiés, mais vise à numériser en priorité  les grands textes fondamentaux et les sources doctrinales de l'histoire du droit.

L'objectif de la liste Pfister-Roumy est de sélectionner des monographiques et des périodiques issus \enquote{uniquement des collections de Cujas}. Cette sélection porte sur les \enquote{fondamentaux du droit à numériser en priorité}. À son origine, les deux professeurs recensent \enquote{471 titres fondamentaux pour la recherche en histoire du droit}, couvrant une vaste période allant du \enquote{XVe siècle au XXe siècle, soit 1462 volumes et 1 001 559 pages}. A cet ensemble patrimonial, les professeurs adjoignent deux périodiques de références: \enquote{Duvergier et le Bulletin des arrêts de la Cour de cassation (Chambre criminelle et Chambres civiles) jusqu'en 1950.}. Ce projet est initialement pensé lors de la candidature de Cujas pour devenir pôle de réplication de Persée, se caractérise par des ambitions techniques élevées, reposant sur une \enquote{numérisation automatique non descriptive, en mode image océrisé, avec une recherche plein texte et la structuration documentaire à granularité fine, emploi de la DTD TEI, réalisé exclussivement en interne}. 

L'abandon de ce projet de réplication oblige la bibliothèque Cujas à repenser son programme de numérisation, contraignant l'établissement à revoir à la baisse le nombre d'ouvrages sélectionnés. Le corpus initial est ainsi réduit, passant de 471 ouvrages à 436 ouvrages. En effet, à la même époque, la BNF développe en parallèle un projet scientifique similaire, intitulé la liste Kerbrat. Élaborée sous la direction du professeur Yann Kerbrat et d'un groupe de quatorze universitaire, cette sélection a pour objectif \enquote{de créer une liste de références à partir des collections de la BnF} dans le domaine du droit. Dans un soucis d'optimisation des moyens financiers et de cohérence scientifique, une concertation étroite s'établit entre les deux institutions afin de coordonner leurs campagnes respectives et d'éviter les doublons. Ainsi, lorsque la bibliothèque Cujas conserve un ouvrage identfique à celui de la BNF, il lui est néanmoins permis de le numériser, mais à la condition expression de sélectionner une édition différente de celle de la BNF. Certains titres, impossibles à substituer par une autre édition, sont alors purement et simplement retirés de la liste de Cujas. 

Au-delà des arbitrages scientifiques, la mise en oeuvre de la liste Pfister-Roumy se heurte rapidement aux contraintes matérielles du scanner acquis par la bibliothèque Cujas. Cet équipement s'avère en effet beaucoup moins efficace que prévu, confrontant les agents à de \enquote{nombreux problèmes liés à son utilisation : il n'est pas possible d'utiliser toutes les fonctions automatiques de numérisation du scanner, forte probabilité de devoir passer en manuel pour les papiers anciens, difficultés de réglage et fréquents dysfonctionnements nécessitant d'éteindre et de rallumer les appareils et de refaire tous les réglages à chaques fois}. De surcroît, le scanner fait peser un fort risque de dégradation physique sur les documents originaux, dont la fragilité exige d'extrêmes précautions. 

C'est la raison pour laquelle la bibliothèque Cujas choisit de déléguer la numérisation des volumes les plus vulnérables de sa liste à un prestataire externe. Malgré les difficultés logistiques pressenties face à l'ampleur d'un tel projet, l'établissement valide la liste Pfister-Roumy en articulant sa production autour d'une chaîne hybride. D'une part, Cujas réemploie sa chaîne de numérisation interne, développée et stabilisée lors de l'exposition de Jean Carbonnier; d'autre part, elle externalise le traitement des ouvrages antérieurs à 1830 afin de garantir leur préservation. Cette répartition du travail permet d'alimenter l'ambitieux projet de créer une bibliothèque numérique juridique fédérée unie par le protocole de moissonnage OAI-PMH. Cette liste Pfister-Roumy constitue ainsi le véritable pilier fondateur de la bibliothèque numérique CujasNum.